% ============================================================
% 第2章：系统方案与架构设计
% ============================================================

\section{总体架构设计}
\label{sec:architectrue}
系统采用“服务层—流水线推理层—工具层—数据层”的四层架构，将自然语言需求输入到三类工程文档输出的全过程组织为可自动执行的处理链路。各层之间职责边界清晰，数据流转方向限定为自上而下的逐层加工过程，从而提升系统的模块化程度、可维护性和可扩展性。

服务层基于 FastAPI 构建，对外提供五个核心端点：\texttt{GET /health}（健康检查）、\texttt{POST /analyze}（返回 \texttt{SelectionReport} JSON）、\texttt{POST /replacement}（返回 \texttt{ReplacementReport} JSON）、\texttt{POST /agent/chat}（支持 \texttt{session\_id} 多轮对话）和 \texttt{GET /agent/sessions}（活跃会话列表），并统一处理请求路由、参数校验与 UTF-8 中文编码兼容。

流水线推理层实现了"需求解析—器件检索—混合评分—证据构建—报告生成"五阶段线性处理流程：自然语言需求经四级容错解析链转换为结构化约束对象，随后通过 eZ-PLM API 多前缀分组查询获取候选器件，再经双模自适应评分引擎排序，并为每个器件生成证据链条目，最终输出 BOM 选型清单、风险评估报告和拓扑分析报告三类工程文档。各阶段的详细实现见第四章。

工具层封装了六类可独立测试和替换的核心能力模块：eZ-PLM API 客户端（HMAC-SHA256 签名与响应解析）、离线器件数据管理器、混合评分引擎（规则/LLM 双模切换）、证据构建器（八类字段级证据生成）、报告生成器（BOM/风险评估/拓扑分析）和替代器件查找器（国产替代与功能等效替代搜索）。

数据层汇聚三类异构数据源：eZ-PLM API（物料查询与参考设计查询接口）、本地器件数据库（JSON 格式，覆盖 Buck/Boost/LDO）和 RAG 工程知识库（ChromaDB 向量存储，29 条条目覆盖 11 个专业类别，通过 \texttt{all-MiniLM-L6-v2} 嵌入为 384 维向量）。其中，离线数据库在 API 不可用时作为兜底数据源；RAG 检索结果在每次 Pipeline 分析时自动注入三类输出报告。

表\ref{tab:tech}汇总了系统核心技术栈的选型方案与对应的依据。
\begin{table}[htbp]
\centering
\caption{核心技术栈选型及依据}
\label{tab:tech}
\footnotesize
\setlength{\tabcolsep}{4pt}
\renewcommand{\arraystretch}{0.95}
\begin{tabularx}{\textwidth}{p{2cm}p{4cm}X}
\toprule
\textbf{类别} & \textbf{选用方案} & \textbf{选用依据} \\
\midrule
大语言模型 & DeepSeek & 支持高精度中文语义理解，API 调用成本可控，统一接口可切换至其他 OpenAI 兼容模型 \\
后端框架 & FastAPI 0.136 + Uvicorn & 异步高性能，自动生成 OpenAPI 文档，深度集成 Pydantic v2，支持类型安全校验 \\
数据建模 & Pydantic v2（2.13） & 声明式字段校验，支持 JSON 自动序列化，便于数据一致性验证 \\
向量数据库 & ChromaDB 1.5 & 支持持久化增量写入，余弦相似度检索满足工程知识语义匹配需求 \\
嵌入模型 & all-MiniLM-L6-v2 & 模型体积约 80\,MB，384 维向量兼顾检索速度与精度，适用于中英文混合语义检索 \\
Agent 框架 & LangChain 1.3 & 提供工具调用、会话记忆与多轮推理能力，ReAct 推理链可用 \\
前端展示 & Streamlit 1.24+ & 纯 Python 可视化交互界面，支持实时渲染评分仪表盘及证据链对比表 \\
\bottomrule
\end{tabularx}
\end{table}

\section{评分模型设计}
\subsection{双模自适应评分框架}

为兼顾评分稳定性与智能化程度，系统设计了双模自适应评分框架：仅当 LLM 服务可用且候选器件具备 eZ-PLM 参考设计数据时，启用 LLM 增强模式；否则自动回退至纯规则模式，保证外部模型或数据不可用时系统仍可稳定运行。

\subsection{纯规则模式}

当 LLM 不可用或缺少参考设计数据时，系统采用基于参数匹配和供应链稳定性的双维评分模型：

\begin{equation}
S_{\text{total}} = 0.80 \times S_{\text{param}} + 0.20 \times S_{\text{supply}}
\label{eq:rule}
\end{equation}

上述权重设定遵循电子元器件选型领域的多属性决策方法论：电气参数匹配是决定器件可用性的首要硬约束，故参数匹配权重占主导（0.80）\cite{pecht2008parts}；供应链稳定性权重（0.20）则参照库存水平与生命周期状态两大类指标的综合赋权研究\cite{chen2021supplyrisk,goel2021supply}。

参数匹配分 $S_{\text{param}}$ 采用动态约束归一化方法计算：

\begin{equation}
S_{\text{param}} = \frac{\displaystyle\sum_{i=1}^{N} s_i}{N} \times 100
\label{eq:param}
\end{equation}

式中，$N$ 表示实际存在的参数约束项数量，$s_i \in [0,1]$ 表示对应约束项的归一化得分。当用户未提供明确参数约束时，系统将参数匹配分设为中性值 50，以避免对候选器件产生系统性惩罚。

\subsection{LLM 增强模式}

当 LLM 服务可用且候选器件具备参考设计数据时，评分模型由双维规则评分扩展为四维混合评分：

\begin{equation}
S_{\text{total}} = 0.40 S_{\text{param}} + 0.10 S_{\text{supply}} + 0.25 S_{\text{app}} + 0.25 S_{\text{design}}
\label{eq:hybrid}
\end{equation}

其中 $S_{\text{app}}$ 和 $S_{\text{design}}$ 为基于参考设计数据的 LLM 软评分维度（应用场景适配度、设计成熟度），参数匹配和供应链权重分别调整为 0.40 和 0.10，释放的 0.50 权重分配至两个 LLM 维度，遵循“硬约束主导、软知识增强”的分层评估思路\cite{wang2022selection}。

为降低 LLM 输出不确定性对评分结果的影响，系统将 $S_{\text{app}}$ 和 $S_{\text{design}}$ 的取值范围限定为 $[0,100]$，并对模型返回结果进行格式校验与数值边界检查。若返回结果格式不完整、字段缺失或分值越界，系统将触发回退机制，改用纯规则评分结果。

\subsection{去重、排序与推荐分级}
eZ-PLM API 返回同一型号族时可能包含多个仅在封装后缀或卷带标识上存在差异的变体（如 LM2596S-5.0/NOPB、LM2596SX-5.0/NOPB）。这些器件核心电气特性一致，若全部纳入推荐列表易造成信息冗余。

系统通过 \texttt{\_extract\_core\_mpn()} 提取核心 MPN——去除 “/NOPB” 等质量等级后缀、“-TR” 等卷带标识及封装变体字符——以核心 MPN 为键分组去重，每组保留综合得分最高者，并将被合并变体记录至 \texttt{reasons} 字段以保留可追溯性。

去重完成后，系统按照综合得分降序排序，并采用 Top-N 分级策略生成推荐结果。综合得分不低于 75 分且排名前五的器件被标注为“推荐”（recommended）；排名第六至第十五且综合得分不低于 50 分的器件被标注为“备选”（backup）；其余器件则标注为“不推荐”（not\_recommended）。最终输出结果最多保留 15 条记录，以兼顾推荐结果的信息完整性和前端展示的可读性。
\section{数据模型设计}

为了同时满足 eZ-PLM 平台的数据互通需求与三类工程文档的结构化输出需求，系统以 Pydantic 为基础定义了一套标准化的 IR 数据模型体系, 该体系由以下五个核心数据模型构成：

（一）\texttt{PartIR} 由 21 个字段组成，完整描述元器件的可检索属性，见表\ref{tab:partir}。
\begin{table}[h]
\centering
\caption{\texttt{PartIR} 字段说明}
\label{tab:partir}
\footnotesize
\setlength{\tabcolsep}{4pt}
\renewcommand{\arraystretch}{0.9}
\begin{tabularx}{\textwidth}{p{2cm}p{4cm}X}
\toprule
\textbf{类别} & \textbf{字段名} & \textbf{说明} \\
\midrule

\multirow{2}{*}{制造商信息}
& \texttt{part\_number} & 制造商型号 (MPN)\\
& \texttt{manufacturer} & 制造商名称 \\
\midrule

\multirow{4}{*}{电气参数}
& \texttt{input\_voltage\_min\_v} & 最小输入电压（V） \\
& \texttt{input\_voltage\_max\_v} & 最大输入电压（V） \\
& \texttt{output\_voltage\_v} & 固定输出电压（V） \\
& \texttt{output\_current\_max\_a} & 最大输出电流（A） \\
\midrule

\multirow{2}{*}{物理属性}
& \texttt{package} & 封装类型 \\
& \texttt{temperature\_min\_c} / \texttt{temperature\_max\_c} & 工作温度范围（°C） \\
\midrule

分类标识
& \texttt{category} / \texttt{topology} & 器件类别与拓扑类型 \\
\midrule

国产标识
& \texttt{is\_domestic} & 是否国产器件 \\
\midrule

\multirow{2}{*}{认证与合规}
& \texttt{automotive\_grade} & 是否通过 AEC-Q100 认证 \\
& \texttt{lifecycle\_status} & 生命周期状态 \\
\midrule

\multirow{3}{*}{供应链状态}
& \texttt{stock} & 库存数量 \\
& \texttt{unit\_price\_cny} & 参考单价（人民币） \\
& \texttt{replacement\_for} & 可替代型号列表 \\
\midrule

\multirow{3}{*}{数据溯源}
& \texttt{datasheet\_url} & 数据手册链接 \\
& \texttt{ezplm\_part\_id} & eZ-PLM 标识符 \\
& \texttt{source} & 数据来源 \\
\midrule

描述
& \texttt{description} & 功能描述 \\
\bottomrule
\end{tabularx}
\end{table}

（二）\texttt{ScoreBreakdown} 记录候选器件的维度得分与计算依据，包含子维度分数（\texttt{parameter\_match\_score}、\texttt{supply\_risk\_score}、\texttt{total\_score}）、评分原因列表（\texttt{reasons}）、评分模式标识（\texttt{rule\_only} 或 \texttt{llm\_enhanced}），以及 LLM 子评分结果，保证评分结论可追溯至具体计算依据。

(三) \texttt{TopologyIR} 采用四子模型嵌套结构，字段定义见表\ref{tab:topoir}。
\begin{table}[h]
\centering
\caption{\texttt{TopologyIR} 字段说明}
\label{tab:topoir}
\footnotesize
\setlength{\tabcolsep}{4pt}
\renewcommand{\arraystretch}{0.9}
\begin{tabularx}{\textwidth}{p{3.2cm}p{4.5cm}X}
\toprule
\textbf{子模型} & \textbf{字段名} & \textbf{说明} \\
\midrule
\multirow{6}{*}{\texttt{TopoNode}}
& \texttt{node\_id} & 节点唯一标识 \\
& \texttt{node\_type} & 节点类型 \\
& \texttt{label} & 中文标签 \\
& \texttt{description} & 功能描述 \\
& \texttt{connected\_to} & 邻接节点 ID 列表 \\
& \texttt{signal\_type} & 信号类型 \\
\midrule

\multirow{8}{*}{\texttt{ExternalComponent}}
& \texttt{refdes} & 参考位号 \\
& \texttt{component\_type} & 元件类型 \\
& \texttt{value} & 推荐参数值 \\
& \texttt{tolerance} & 容差 \\
& \texttt{voltage\_rating} & 耐压等级 \\
& \texttt{temperature\_coeff} & 温度系数 \\
& \texttt{package} & 封装规格 \\
& \texttt{selection\_notes} & 选型备注 \\
\midrule

\multirow{7}{*}{\texttt{ThermalEstimate}}
& \texttt{estimated\_efficiency\_pct} & 估算效率（\%） \\
& \texttt{output\_power\_w} & 输出功率（W） \\
& \texttt{total\_power\_loss\_w} & 总功耗（W） \\
& \texttt{junction\_temp\_c} & 结温估算值（°C） \\
& \texttt{junction\_temp\_max\_c} & 最高允许结温（°C） \\
& \texttt{thermal\_margin\_c} & 热余量（°C） \\
& \texttt{needs\_heatsink} & 是否需要散热器 \\
\bottomrule
\end{tabularx}
\end{table}

（四）\texttt{RiskIR} 描述选型方案的整体风险等级（低/中/高）及结构化风险条目。每条 \texttt{RiskItem} 均由规则引擎生成，包含风险类型、严重程度、风险描述、相关器件型号、缓解建议和证据引用字段。规则引擎已覆盖 13 类确定性风险场景（无候选器件、库存不足、停产器件、供应商过度集中、关键字段缺失等），LLM 仅负责补充工程背景说明，不参与风险判定或新增风险条目，保证风险结论建立在可验证的事实基础之上。


（五）\texttt{SelectionReport} 和 \texttt{ReplacementReport} 分别为选型分析和替代搜索的顶层输出结构，封装了用户输入、结构约束、候选/推荐器件列表、风险汇总、证据链和摘要文本，为国产替代与功能等效替代两类场景提供统一的数据返回格式。
